\documentclass[a4paper,11pt]{article}
\usepackage[utf8]{inputenc}
\usepackage[T1]{fontenc}
\usepackage[french]{babel}
\usepackage[left=2cm,right=2cm,top=2cm,bottom=2cm]{geometry}
\usepackage{xcolor}
\usepackage{hyperref}
\usepackage{fontawesome5}
\usepackage{parskip}

% --- Couleurs LIC ---
\definecolor{primary}{RGB}{0, 60, 120}
\hypersetup{colorlinks=true, urlcolor=primary}

\begin{document}

% --- EXPÉDITEUR ---
\noindent
\begin{minipage}[t]{0.5\textwidth}
    \textbf{Loïc Aron MBASSI EWOLO}\\
    Élève-Ingénieur Bac+5 -- Double diplôme ENSEA\\[3pt]
    \faMapMarker\ Cergy, France\\
    \faPhone\ (+33) 7 59 52 33 98\\
    \faEnvelope\ \href{mailto:loicaron.mbassiewolo@ensea.fr}{loicaron.mbassiewolo@ensea.fr}
\end{minipage}
\hfill
% --- DATE ---
\begin{minipage}[t]{0.4\textwidth}
    \raggedleft
    Cergy, le 10 mars 2026
\end{minipage}

\vspace{1.2cm}

% --- DESTINATAIRE ---
\hfill
\begin{minipage}[t]{0.5\textwidth}
    \raggedleft
    \textbf{LIC -- L'Informatique Communicante}\\
    M. Eric NIZARD, Directeur\\
    2 rue Maurice Hartmann\\
    92130 Issy-les-Moulineaux
\end{minipage}

\vspace{1cm}

\noindent\textbf{Objet :} Candidature en alternance -- Ingénierie des objets connectés sécurisés (Réf. 2601IOCS)

\vspace{0.8cm}

Madame, Monsieur,

L'analyse de protocoles de communication sécurisés (ISO 7816, NFC, BLE) appliquée aux terminaux de paiement et objets connectés constitue un domaine dans lequel je souhaite approfondir mon expertise. L'offre d'alternance proposée par LIC, au croisement de la sécurité numérique, de l'électronique et des protocoles de paiement, correspond à mes compétences et à mon parcours d'ingénieur en systèmes embarqués.

En double diplôme à l'ENSEA (systèmes embarqués, signal, IA) après un cursus en informatique à l'École Polytechnique de Yaoundé, j'ai acquis une base solide en programmation bas niveau (C/C++, STM32, Arduino) et en analyse de systèmes. Au laboratoire IOTA de l'ENSPY, j'ai participé à la conception de gants IoT instrumentés (capteurs IMU, flex sensors, communication sans fil), ce qui m'a confronté aux enjeux d'intégration hardware/software et de fiabilité des communications entre objets connectés. Mon stage au IoT Lab (TinyML, TF Lite sur microcontrôleurs) a consolidé cette expertise en systèmes contraints.

Mon stage de recherche à INRIA Saclay (équipe TRIBE, avril-août 2026) porte sur l'analyse automatisée de la sécurité et de la conformité d'applications mobiles : analyse statique de SDKs, évaluation des permissions et de la transparence des politiques de confidentialité (NLP, RGPD). Ce travail m'a permis de développer une rigueur dans l'analyse de conformité et la rédaction de spécifications techniques, compétences directement transposables à l'analyse de protocoles et à la qualification d'outils de tests chez LIC. Par ailleurs, mon expérience chez CSC (secteur Fintech) m'a familiarisé avec les normes de sécurité financière et le développement de plateformes de paiement, un contexte proche de celui de vos clients (EMVCo, MasterCard).

Je pratique l'anglais au quotidien dans mes travaux de recherche (rédaction scientifique, échanges avec des partenaires internationaux, stage au MILA à Montréal). Je suis disponible dès septembre 2026 pour une alternance de 12 mois, à l'issue de mon stage INRIA.

Je reste à votre disposition pour un entretien afin de discuter de ma candidature et de la valeur que je peux apporter à vos projets de certification et de conseil.

Je vous prie d'agréer, Madame, Monsieur, l'expression de mes salutations distinguées.

\vspace{0.8cm}

\textbf{Loïc Aron MBASSI EWOLO}\\
\textit{Élève-Ingénieur ENSEA / Polytechnique Yaoundé}

\end{document}
